\originalsection{旧期中硬球减速与旧期末旋转箍缩}
\begin{exercise}\label{ex:hist05-mid-2}
2005期中原题2。半径均为$R$的两硬球以碰撞参数$b$相遇，球1初速$V_0\hat x$，球2静止。题面先给球2质量无穷时的球1末速度：$b<2R$下
\begin{equation}\mathbf V_{\infty}'=V_0\left[\left(\frac{b^2}{2R^2}-1\right)\hat x+\frac bR\sqrt{1-\frac{b^2}{4R^2}}\hat y\right],\end{equation}
$b>2R$时仍为$V_0\hat x$。(a)现球1、球2质量分别$m_1,m_2$，求球1末速度；(b)球1束穿过静止球2密度$n$的介质，用$\dd\mathbf V_0/\dd t=-\nu_{\rm sd}\mathbf V_0$定义减速率，求$\nu_{\rm sd}$。
\end{exercise}
\begin{solution}
编者独立解答。令$M=m_1+m_2$，碰撞法线几何决定相对速度转角$\chi$，$\cos\chi=b^2/(2R^2)-1$、$\sin\chi=(b/R)\sqrt{1-b^2/(4R^2)}$。质心速度$m_1V_0\hat x/M$不变，故$b<2R$时
\begin{equation}
\mathbf V_1'=\frac{m_1V_0}{M}\hat x+\frac{m_2V_0}{M}(\cos\chi\hat x+\sin\chi\hat y),
\quad
\mathbf V_2'=\frac{m_1V_0}{M}[\hat x-(\cos\chi\hat x+\sin\chi\hat y)].
\end{equation}
$b\ge2R$不碰撞。横向平均相消，单位时间碰撞参数环面积是$nV_0(2\pi b\dd b)$，从而
\begin{equation}
\nu_{\rm sd}=nV_0\frac{m_2}{M}\int_0^{2R}(1-\cos\chi)2\pi b\dd b
=4\pi R^2 nV_0\frac{m_2}{m_1+m_2}.
\end{equation}
这定义的是束平均定向动量衰减，不是把单次球1速度大小按相同系数衰减；$m_2\to\infty$回到题给硬壁结果。
\end{solution}

\begin{exercise}\label{ex:hist05-final-3}
2005期末原题3。轴对称$\theta$箍缩刚体转动$\mathbf v=r\Omega\hat\theta$，中心$\beta_0=p_0/(B_0^2/2\mu_0)=0.5$，边缘$a\Omega=\sqrt{2k_BT_0/m_i}$；$T_e=T_i$，忽略黏性，两小问独立。(a)若$p(a)=0$，求边缘内$B(a^-)/B_0$；(b)另设$n=n_0$、$T(r)=T_0(1-r^2/a^2)^2$，求单位长度热能$W$、传导功率$P$及$\tau_E=W/P$，以$a$、中心碰撞时间$\tau_0$、中心热Larmor半径$\rho_L$和$B(a^-)/B_0$表示，并说明粒子种类。
\end{exercise}
\begin{solution}
编者独立解答。轴向$B=B_z(r)$的径向力平衡是
\begin{equation}\frac{\dd}{\dd r}\left(p+\frac{B^2}{2\mu_0}\right)=m_in(r)\Omega^2r.\end{equation}
故(a)的无附加剖面严格结果为$B(a^-)^2=B_0^2+2\mu_0p_0+2\mu_0m_i\Omega^2\int_0^a n(r)r\dd r$。原题(a)只给中心$p_0,T_0$而未给$n(r)$，因此边缘磁场数值并不唯一；若另附加均匀密度，$n_0k_BT_0=p_0/2$及$a^2\Omega^2=2k_BT_0/m_i$给$B(a^-)^2/B_0^2=1+\beta_0+\beta_0/2=1.75$。不能将(b)的均匀密度暗中作为(a)原有条件。

在(b)中$W=\int_0^a(3/2)(nT_e+nT_i)2\pi r\dd r=\pi n_0 k_BT_0a^2$，这里$T_0$为K。若选强磁化离子经典闭合$\kappa_{i\perp}=C_i n_0k_B\rho_{Li}^2/\tau_{ii}$，定义$\rho_{Li}^2=2k_BT_0m_i/(e^2B_0^2)$，且$\tau_{ii}(r)=\tau_{ii,0}[T(r)/T_0]^{3/2}$，则$\kappa_{i\perp}(r)=\kappa_0(B_0/B(r))^2[T(r)/T_0]^{-1/2}$，$\kappa_0=C_i n_0k_B\rho_{Li}^2/\tau_{ii,0}$。虽然$T(a)=0$，$\kappa\partial_rT$极限有限：$q_a=4\kappa_0T_0(B_0/B_a)^2/a$，于是
\begin{equation}
P=8\pi\kappa_0T_0(B_0/B_a)^2,\qquad
\tau_E=\frac{a^2\tau_{ii,0}}{8C_i\rho_{Li}^2}\left(\frac{B_a}{B_0}\right)^2.
\end{equation}
等温等密度下离子垂直导热通常支配该经典估计，所以碰撞时间和热回旋半径取离子—离子量。$C_i$取决于采用的碰撞时间／输运系数定义；题PDF声称“下一页附Formulary表”，实际下一页已是原题4，附表未在这份八页PDF中取得。故本稿把这一未给的无量纲系数保留，不冒充来源给出唯一数值前因子；$T(a)=0$也超出弱耦合强磁化闭合的严格有效域。
\end{solution}
